\section{Related work}\label{sec:literature}

Two-stage capacitated facility location entered the operations research
literature as a distribution-planning model. General surveys of facility
location and of multi-level location are given by \citet{KloseDrexl2005}
and \citet{OrtizAstorquizaEtAl2018}.
\citet{PirkulJayaraman1998} formulated the multiproduct two-stage
capacitated version, observed its \NP-hardness and solved it by
Lagrangian relaxation, and this is the formulation that the later
benchmark work inherits. Exact and heuristic effort has continued on both
the single- and multi-stage variants. \citet{LitvinchevOzuna2012} derived
Lagrangian bounds for two-stage capacitated location, together with an
accompanying heuristic, and \citet{FernandesEtAl2014} proposed a simple
and effective genetic algorithm for the single-product two-stage problem
along with the large benchmark instance sets that the multiproduct
studies later reused. On the exact side, Benders decompositions tailored
to facility location are a mature computational line
\citep{FischettiLjubicSinnl2017}, and the disaggregated cuts of
Section~\ref{sec:certificates}, derived from the routing oracle of
Proposition~\ref{prop:oracle}, sit in that tradition. The benchmark this paper builds on is due
to \citet{MauriEtAl2021}, who assembled a hundred multiproduct instances
and, finding that a commercial mixed-integer solver closed none of them
within the imposed limits, proposed two hybrid metaheuristics, a
clustering search matheuristic that calls CPLEX on network-flow
subproblems and a biased random-key genetic algorithm with local search.
\citet{MoraisEtAl2022} then addressed the same multiproduct problem with
a hybrid biased random-key genetic algorithm that adds local branching,
extending a single-product hybrid from the same group. Across this line the model is treated as empirically hard, and
the hardness is managed rather than explained.

The complexity of facility location itself is well understood in the
uncapacitated single-product case, and it is where our covering arguments
originate. Without the metric assumption, uncapacitated facility location
contains Set Cover, so \citet{Hochbaum1982} obtained a greedy
$O(\log n)$ approximation, and the lower bounds recorded next show that no
better guarantee is possible in general unless $\Ppoly=\NP$.
The matching lower bounds are those of \citet{Feige1998} and, under the
weaker assumption \Ppoly${}\neq{}$\NP, of \citet{DinurSteurer2014},
which together pin the Set Cover threshold at $(1-\varepsilon)\ln n$.
These bounds transfer to the depot-to-customer stage of MP-TSCFLP and
account for its inapproximability. The classical hardness of the
underlying decision problems traces back to \citet{Karp1972}, and the
distinction between weak and strong \NP-hardness that organises our
numeric layer is that of \citet{GareyJohnson1979}.

Parameterized complexity has been brought to bear on optimisation and
covering problems for two decades, and the vocabulary and lower-bound
machinery we use are standard \citep{DowneyFellows2013,Cygan2015}. The
$\mathrm{W}[2]$-hardness of \textsc{Set Cover} parameterized by the
solution size is classical \citep{DowneyFellows2013}, the ETH-based lower
bounds we compose with are due to \citet{ChenHuangKanjXia2006}, and the
compression arguments in the kernelization section rely on
the number-reduction technique of \citet{FrankTardos1987}. The parameterized
view of location and network problems has grown quickly, though along axes
different from ours. The closest antecedent is the parameterized view of
facility location by \citet{FellowsFernau2011}, published in this
journal, which establishes W-hardness and \FPT\ results for
uncapacitated facility location variants under related parameters.
Capacities, a second stage and multiple products all lie outside its
scope, and that is the gap the present analysis fills.
Facility-location-flavoured clustering admits tight
fixed-parameter approximations in the number of centres
\citep{CohenAddadEtAl2019}. Under uniform capacities tight \FPT\
approximations in the number of centres remain available, and
constant-factor \FPT\ approximations with larger factors extend to the
general capacitated setting \citep{CohenAddadLi2019}. The survivable
network design problem is fixed-parameter tractable in its edge- and
element-connectivity forms while turning $\mathrm{W}[1]$-hard for vertex
connectivity \citep{FeldmannMukherjeevanLeeuwen2025}. Covering objectives
on structured graphs, tractable by treewidth alone in their nonnegative
form, in general require treewidth combined with maximum degree
\citep{FominRamamoorthi2022}. Each of the clustering, connectivity and
covering results turns on a metric, a connectivity
requirement or a graph width, whereas the hardness we isolate is arithmetic,
carried by capacities and demands, so their positive machinery does not
transfer to the representations analysed here, while richer
representations remain unanalysed rather than excluded. The contrast is sharpest on the
approximation side. What separates our covering core from the capacitated clustering results
is the absence of a metric. There the parameterized lower bounds have
recently caught up with the classical ones, since approximating
\textsc{Set Cover} parameterized by the solution size within any constant
factor is itself $\mathrm{W}[2]$-hard \citep{LinRenSunWang2023}. This suggests the same barrier for
MP-TSCFLP, a transfer we state as a conjecture among the open problems
of Section~\ref{sec:conclusion}. Closest in
spirit to the present analysis is the parameterized study of
minimum-weight solution subgraphs of and-or graphs by
\citet{DantasProttiSouza2013}, where a cheapest subnetwork must respect
local dependency semantics and zero-weight arcs are exactly what turns a
fixed-parameter tractable problem $\mathrm{W}[2]$-hard. Zero fixed charges play a
comparable role in our constructions, which require either large numbers,
many products or positive charges to produce hardness.

Table~\ref{tab:positioning} closes the section by placing this paper
against the works discussed above, one line per selected work and one
column per feature, with a mark recording only what the corresponding
paragraph of this section states. The last line of the table records a
combination of features that no previous line fills, a parameterized
analysis of the two-stage capacitated multiproduct model.

\begin{table}[H]
\footnotesize
\centering
\caption{Positioning against the related work of this section. A check
mark means the feature is treated in the cited work as described above.
The circle marks bounds or partial guarantees claimed without an
exactness or completeness proof. Blank means not
treated.}
\label{tab:positioning}
\resizebox{\textwidth}{!}{%
\begin{tabular}{lcccccc}
\toprule
 & two-stage & multi- & capa- & exact or & heuris- & parameterized \\
 & location & product & cities & lower bounds & tics & analysis \\
\midrule
\citet{PirkulJayaraman1998} & \checkmark & \checkmark & \checkmark &
$\circ$ & \checkmark & \\
\citet{LitvinchevOzuna2012} & \checkmark & & \checkmark & $\circ$ &
\checkmark & \\
\citet{FernandesEtAl2014} & \checkmark & & \checkmark &
 & \checkmark & \\
\citet{MauriEtAl2021} & \checkmark & \checkmark & \checkmark & &
\checkmark & \\
\citet{MoraisEtAl2022} & \checkmark & \checkmark & \checkmark & &
\checkmark & \\
\citet{DantasProttiSouza2013} & & & & \checkmark & & \checkmark \\
\citet{CohenAddadLi2019} & & & \checkmark & $\circ$ & &
\checkmark \\
\citet{FeldmannMukherjeevanLeeuwen2025} & & & & \checkmark & &
\checkmark \\
This paper & \checkmark & \checkmark & \checkmark & \checkmark & &
\checkmark \\
\bottomrule
\end{tabular}}
\end{table}
